% Section 3 — What we found (double-width box: big figure left, text/eqn right).
\begin{block}{3.\quad Ghost-free amplification and suppression}
  \begin{columns}[c]
    % --- big corner figure ---
    \begin{column}{0.60\linewidth}
      \centering
      \includegraphics[width=\linewidth]{result_landscape.pdf}
    \end{column}
    % --- interpretation + dominant operator ---
    \begin{column}{0.38\linewidth}
      \raggedright
      The conversion responds to the full content of
      each theory, but a few operators are \emph{robustly dominant}. The most
      consistent amplifier is
      \vspace{0.8ex}

      % Dominant-operator equation inside a padded coloured box (a plain
      % beamercolorbox, so the display equation cannot escape below the fill).
      {\centering
        \begin{beamercolorbox}[wd=\linewidth,sep=1.6ex,center]{eqn box}
          % \resizebox scales the equation to the box width on a SINGLE line --
          % maximally large for emphasis, and guaranteed never to wrap (the column
          % is too narrow for a fixed \Huge/\huge, which broke after \supset).
          \resizebox{0.92\linewidth}{!}{%
            $\lag \;\supset\; \rcD{^{\alpha}}\MAGT{^{\beta\gamma\delta}}\,
              \MAGF{_{\alpha\beta\gamma\delta}}$}
        \end{beamercolorbox}\par}
      \vspace{0.8ex}
      Some directions in the coupling space \textbf{amplify} the conversion and
      others \textbf{suppress} it --- distinct regions, not a single effect. The
      amplification persists without
      propagating-torsion, so purely ghost-driven growth cannot explain it.
    \end{column}
  \end{columns}
\end{block}
